% Part: computability
% Chapter: recursive-functions
% Section: primes

\documentclass[../../../include/open-logic-section]{subfiles}

\begin{document}

\olfileid{cmp}{rec}{pri}
\olsection{质数}

有界量词和有界极小化为我们提供了大量工具，用于证明常见的函数与关系是原始递归的。例如，考虑关系“$x$ 整除 $y$”，记作 $x \mid y$。关系 $x \mid y$ 成立，当且仅当 $y$ 除以 $x$ 没有余数，也就是说，$y$ 是 $x$ 的整数倍。（如果不成立，即 $y$ 除以 $x$ 的余数 $> 0$，我们记作 $x \nmid y$。）换言之，$x \mid y$ 当且仅当存在某个 $z$ 使得 $x \cdot z = y$。显然，任何这样的 $z$，如果存在，必定 $\le y$。因此，$x \mid y$ 当且仅当存在某个 $z \le y$ 使得 $x \cdot z = y$。我们可以通过对 $=$ 和乘法进行有界存在量化来定义关系 $x \mid y$：
\[
x \mid y \defiff \bexists{z \leq y}{(x \cdot z) = y}.
\]
由此便证明了 $x \mid y$ 是原始递归的。

一个自然数 $x$ 是\emph{质数}，如果它既不是 $0$ 也不是 $1$，并且只能被 $1$ 和它自身整除。换言之，质数满足：只要 $y \mid x$，则要么 $y = 1$，要么 $y=x$。要检验 $x$ 是否为质数，我们只需对所有 $y \le x$ 检查 $y \mid x$ 是否成立，因为如果 $y > x$，则自动有 $y \nmid x$。因此，关系 $\fn{Prime}(x)$（它在且仅在 $x$ 为质数时成立）可以定义为
\[
\fn{Prime}(x) \defiff x \geq 2 \land \bforall{y \leq x}{(y \mid x \lif y
  = 1 \lor y = x)}
\]
因而也是原始递归的。

质数有 $2$, $3$, $5$, $7$, $11$ 等。考虑函数 $p(x)$，它返回该序列中第 $x$ 个质数，即 $p(0) = 2$, $p(1) = 3$, $p(2) = 5$ 等。（为方便起见，我们常将 $p(x)$ 写作 $p_x$（$p_0=2$, $p_1=3$ 等）。）

如果我们有一个函数 $\fn{nextPrime(x)}$，它返回大于 $x$ 的第一个质数，那么 $p$ 就可以很容易地用原始递归定义：
\begin{align*}
  p(0) & = 2\\
  p(x+1) & = \fn{nextPrime}(p(x))
\end{align*}
由于 $\fn{nextPrime}(x)$ 是满足 $y > x$ 且 $y$ 是质数的最小 $y$，它可以容易地通过无界搜索来计算。但根据 Euclid（欧几里得）的一个结果，它也可以用有界极小化来定义：在 $x$ 和 $\fact{x}+1$ 之间总存在一个质数。
\[
  \fn{nextPrime(x)} =
  \bmin{y \leq \fact{x}+1}{(y > x \land \fn{Prime}(y))}.
\]
这表明 $\fn{nextPrime}(x)$ 进而 $p(x)$ 不仅是可计算的，而且是原始递归的。

（如果你感到好奇，这里给出 Euclid 定理的一个快速证明。假设 $p_n$ 是 $\le x$ 的最大质数，并考虑所有 $\le x$ 的质数的乘积 $p =
p_0\cdot p_1 \cdot \dots \cdot p_n$。要么 $p+1$ 是质数，要么在 $x$ 和 $p+1$ 之间存在一个质数。为什么呢？假设 $p+1$ 不是质数。那么存在某个质数 $q$ 使得 $q \mid p+1$ 且 $q < p+1$。所有 $\le x$ 的质数都不整除 $p+1$。（根据 $p$ 的定义，每个质数 $p_i \le x$ 都整除 $p$，即余数为 $0$。因此，每个质数 $p_i \le x$ 除 $p+1$ 的余数均为 $1$，所以 $p_i \nmid p+1$。）由此，$q$ 是一个大于 $x$ 且小于 $p+1$ 的质数。又因为 $p \le \fact{x}$，所以存在一个大于 $x$ 且 $\le \fact{x}+1$ 的质数。）

\begin{prob}
使用有界极小化定义整数除法 $d(x, y)$。
\end{prob}

\end{document}